\section{The same unit excludes the full normalized adjoint branch}
\label{sec:kt-adjoint-unit-exclusion}

The ordinary rank-two exclusion also applies when
only its adjoint residual representation is known.
The necessary scalar character can be recovered
from that residual adjoint and then removed.  This
is a new exclusion for the specified adjoint
branch; the general stable rank-three branch
remains outside its scope.

\begin{theorem}[KT-ADJOINT-BRANCH-UNIT-EXCLUSION]
Let $p>3$, let $D<0$ be a fundamental discriminant
prime to $p$, put $E=\mathbf Q(\sqrt D)$ and
$\eta=\chi_D$, and let $e\ge2$ be even with
$(p-1)\nmid2e$.  Let $V$ be a continuous absolutely
irreducible two-dimensional representation of
$G_{\mathbf Q}$ over a finite extension of
$\mathbf Q_p$, with
\[
                  \det V=\eta\varepsilon_p^e.             \tag{AJ.1}
\]
Put $A=\operatorname{Ad}^0(V)=\operatorname{End}^0(V)$,
with the conjugation action.  Suppose that:
\begin{enumerate}
\item $A|_{G_E}$ is unramified at all finite places
      away from $p$;
\item the residual semisimplification of $A|_{G_E}$ is
      $1\oplus\bar\varepsilon_p^e
                  \oplus\bar\varepsilon_p^{-e}$;
\item at a place $v\mid p$ of $E$, $A|_{G_{E_v}}$
      has a stable full flag with ordered inertial
      characters
      $(\varepsilon_p^e,1,\varepsilon_p^{-e})$.
\end{enumerate}
Then $B_{1,\eta\omega^e}$ is not a $p$-adic unit.
The assertion holds at split and inert $p$ and
has no upper bound on $e$.
\end{theorem}

\begin{proof}
We may extend the finite coefficient field whenever
necessary to define a stable lattice or a flag.
Write $\psi=\bar\varepsilon_p^e$ and
$d=\operatorname{ord}(\psi|_{G_E})$.  Since
$p\nmid D$, the fields $E$ and $\mathbf Q(\mu_p)$
are disjoint, and GD gives $d>2$.

\paragraph{The residual two-dimensional representation splits.}
Choose a stable lattice in $V$ and reduce it.
The induced lattice in $\operatorname{End}^0(V)$
has the same residual semisimplification as any
other stable adjoint lattice.  Hypothesis~2
therefore makes the finite image of
$\operatorname{Ad}^0(\bar V)|_{G_E}$ triangularizable:
all its composition factors are one-dimensional.
Conjugation on trace-zero matrices is faithful
on $\mathrm{PGL}_2$ when $p\ne2$.
Thus the finite projective image $P_E$ has a
normal $p$-subgroup $U$, namely the kernel of
the diagonal-character map of that triangular
representation, and $P_E/U$ is abelian of
order prime to $p$.

If $U\ne1$, it has a unique fixed point on
$\mathbf P^1$.  Here is the elementary argument.
A central element of order $p$ in $U$ is
unipotent and has a unique fixed point; centrality
forces all of $U$ to preserve that point.
After putting it at infinity, every element of
the finite $p$-group has trivial diagonal ratio,
so $U$ consists of translations and its common
fixed point is unique.  Normality then forces
$P_E$ to preserve that point.

If $U=1$, $P_E$ is abelian and has an element
of projective order greater than two: otherwise
all its adjoint characters would have order at
most two, contrary to $d>2$.
Such an element is semisimple and its centralizer
in $\mathrm{PGL}_2$ is a torus.  Indeed a commuting
element must preserve its two eigenlines;
exchanging them would invert the projective
element and force order at most two.
Consequently $P_E$ lies in that torus.
In both cases the residual representation on
$G_E$ is reducible.

Its semisimplified character ratio, in a suitable
order, is $\psi$.  Restriction of a semisimple
finite-image representation to a normal subgroup
is semisimple by the socle argument in
KT-ORDINARY-RANK2-RESTRICTION-VARIANT.
Conjugation by $G_{\mathbf Q}$ cannot exchange
these two distinct character lines: their ratio
$\psi$ is conjugation-invariant and differs from
$\psi^{-1}$.  Hence
\[
                  \bar V^{\rm ss}=\lambda\oplus\mu
                                                               \tag{AJ.2}
\]
for global characters, with
$\lambda\mu^{-1}=\bar\eta^t\psi$ for some
$t\in\{0,1\}$.

At complex conjugation, $\det\bar V=-1$ by
(AJ.1), whereas $\psi(c)=1$ because $e$ is
even.  Since both $\lambda(c)$ and $\mu(c)$
are signs, their ratio equals their product.
Thus $t=1$.  Combining the ratio and determinant
in (AJ.1) gives $\mu^2=1$.
Let $\nu:G_{\mathbf Q}\to\{\pm1\}$ be the
canonical characteristic-zero lift of $\mu$.
We have proved
\[
          \bar V^{\rm ss}=
                  \bar\nu\oplus\bar\nu\bar\eta\psi.
                                                               \tag{AJ.3}
\]

\paragraph{The adjoint ordinary flag gives the required quotient.}
Because $p\nmid D$, $E_v/\mathbf Q_p$ is trivial
or unramified quadratic; its inertia is the same
as that of $\mathbf Q_p$.
The highest line $L$ of the flag in hypothesis~3
is the unique inertia eigenline of character
$\varepsilon_p^e$: that character occurs once
among the three distinct composition factors.
It is consequently stable under the full
$G_{\mathbf Q_p}$, which normalizes inertia and
preserves the cyclotomic character.

Take $0\ne N\in L\subset\operatorname{End}^0(V)$.
Conjugation preserves $\det N$, while inertia
scales $N$ by $\varepsilon_p^e$.  The latter
character has infinite order, so $\det N=0$.
The trace is zero as well; thus $N$ is a
rank-one nilpotent and
$\ker N=\operatorname{im}N$ is a line in $V$
stable under $G_{\mathbf Q_p}$.
Write $\delta_+$ for its character and
$\delta_-$ for that of the quotient.
The action on $N$ is $\delta_+/\delta_-$, so
on inertia
\[
 \delta_+/\delta_-=\varepsilon_p^e,
 \qquad \delta_+\delta_-=\varepsilon_p^e,
 \qquad \delta_-^2=1.                                    \tag{AJ.4}
\]

The saturated local line in a stable lattice
shows that $\bar\delta_-|_{I_p}$ is one of
$\bar\nu|_{I_p}$ and $\bar\nu\psi|_{I_p}$.
The second is impossible: its square is
$\psi^2\ne1$ on inertia, whereas (AJ.4) has
square one.  Therefore
$\bar\delta_-|_{I_p}=\bar\nu|_{I_p}$.
Both $\delta_-|_{I_p}$ and $\nu|_{I_p}$ take
values in $\{\pm1\}$, on which reduction is
injective.  They are equal.  It follows that
\[
                            V'=V\otimes\nu              \tag{AJ.5}
\]
has an unramified rank-one quotient at $p$.
This argument does not assume an unproved
converse from adjoint crystallinity: it uses
the exact ordered inertia flag in hypothesis~3.
For classical forms unramified at $p$, their
usual crystalline realization is consistent
with this local conclusion but is not needed
as an extra input.

\paragraph{The same scalar twist removes away-prime inertia.}
For $q\ne p$ and $g\in I_q\cap G_E$,
hypothesis~1 gives $A(g)=1$, so $V(g)=z_g I$.
Its determinant is one by (AJ.1), since both
$\eta$ and $\varepsilon_p$ are trivial on
this inertia subgroup.  Hence $z_g\in\{\pm1\}$.
Formula (AJ.3) shows that its reduction is
$\bar\nu(g)$, and injectivity on signs gives
$z_g=\nu(g)$.  Thus $V'(g)=1$.
The twisted representation is therefore
unramified away from $p$ after restriction
to $G_E$.

Finally (AJ.3) gives
$(\bar V')^{\rm ss}=1\oplus\bar\eta\psi$.
The twist preserves absolute irreducibility.
All three hypotheses of
KT-ORDINARY-RANK2-UNIT-EXCLUSION hold for
$V'$ and $a=e$.  A unit $B_{1,\eta\omega^e}$
would contradict that theorem.
\end{proof}

\paragraph{Application to the source and its boundary.}
For any actual Type~I point whose restricted
rank-three system is of the form
$\operatorname{Ad}^0(V)|_{G_E}$ with (AJ.1),
the source's residual matrix, away-prime
unramifiedness and ordinary flag give exactly
the three hypotheses above.  For
$E=\mathbf Q(\sqrt{-3})$, the lower-weight
classical adjoint subspace constructed in
KT-U3-SYM2-INJECT has determinant normalization
(AJ.1) before taking the adjoint.  Thus the
$a=e$ auxiliary unit excludes every actual
source point in this normalized adjoint
subspace, including residual pairs that
initially carry an unknown quadratic scalar
character.  This is stronger than excluding
only the previously normalized Eisenstein
seeds.

The determinant and adjoint-origin hypotheses
are explicit.  For other auxiliary fields the
same exclusion applies whenever this stated
adjoint origin and determinant have been
identified; no additional spectral construction
over those fields is asserted here.
No assertion identifies every
stable rank-three parameter with this family.
Under the two units with exponents $e$ and
$2e$, actual Type~II points and the specified
adjoint Type~I points are excluded; a
stable Type~I source outside this specified
normalized adjoint family is still a possible
remaining obstruction.
